\rnote{I think much of this stuff would fit better in the introduction.}
\ctodo{I removed the banded fro mthe previous section, can we define it here? It's commented out below.}
% \paragraph{Banded matrices}
% We say a (general) matrix $\bfX$ is $\bands$-banded if for all $i,j \in [n]$, $|i -j| \ge \bands$ implies $\bfX\idx{i}{j} = 0$. While this is off-by-one from the typical definition of bandwidth ($\bfX$ has bandwidth $b-1$), our definition will be useful as it will be natural to match $\bands$-banded matrices with $\minsep$-min-separation. Further, for $\bands$-banded lower-triangular matrices (which will play a central role), $\bands$ intuitively gives the number of bands in the matrix.  

% We may want to move some of this to related work.
While \multiepoch-participation is a good fit for centralized training infrastructure, in cross-device federated learning the order in which data is processed cannot be tightly controlled, as devices make local decisions about if and when they will participate in training \citep{FLO,bonawitz2019towards}. This led to the development of alternative DP training protocols that could benefit from some amplification, such as that of \citet{balle2020privacy}. The cross-device FL setting was in fact a primary motivation for the development of the DP-FTRL algorithm \citep{kairouz2021practical}, which allows performance competitive with DP-SGD analyzed with amplification-via-sampling, but applies to arbitrary data processing patterns (with sensitivity calculated based on the minimum separation between participations).  \citet{denisov22matfact} improved on this approach in the single-participation setting by using matrix factorization in place of tree aggregation; the dense lower-triangular matrix mechanisms introduced in this work could be used in cross-device FL, but only if each client was restricted to participate at most once.  \citet{choquette22multiepoch} extended to \citet{denisov22matfact}, but focused on \multiepoch-participation, which as noted above is only feasible in centralized training scenarios. Hence the key open question we address here: \emph{Can matrix-factorization DP-FTRL (\MFFTRL) be extended to the cross-device federated learning setting with multiple client participations?} 

This question is of substantial practical importance, as \citet{choquette22multiepoch} showed \MFFTRL can substantially outperform tree-aggregation-based DP-FTRL. However, in practice the limitation to a single participation (single epoch) is prohibitive in many private training scenarios, as increasing the number of clients per round (equivalent to the batch size in centralized training) and hence the total epochs of training is perhaps the single most important technique in improving privacy/utility tradeoffs in differentially private machine learning \citep{ponomareva2023dpfy}.

It is also worth clarifying exactly what the extension of \MFFTRL to cross-device FL requires. An arbitrary factorization $\bfA = \bfB \bfC$ could be used under any participation schema (including \bparticipation). However, in order to make deploying such a mechanism useful, it is essential that we can tightly bound $\sensd(\bfC)$, and hence state tight privacy guarantees. Of course, even this is not enough: we could take the trivial factorization $\bfB \leftarrow \bfA$ and $\bfC \leftarrow \bfI$, and tightly bounding sensitivity would be trivial. Hence, we need to be able to both generate ``good'' factorizations (\btodo{Need to reference, possibly forward, a defintion of ``good''}) and bound their sensitivity. The best way to do this is to find a way to efficiently represent the constraint $\sensd(\bfC) \le 1$ in the mathematical program solved to find \emph{optimal} factorizations $\bfA = \bfB \bfC$.

Our approach to extending \MFFTRL to cross-device FL thus begins with a generalization of \multiepoch-participation, \bparticipation, which can practically be enforced by cross-device FL systems \btodo{Will need to say more about this somewhere.}  Then, we introduce a class of matrix mechanisms that perform extremely well while allowing for the efficient computation of sensitivity, and hence, efficient convex optimization to find the necessary matrix factorizations.

We define a participation schema where the distance between any two participations is \emph{at least} \minsep, rather than \emph{exactly} $\minsep$ as in \multiepoch-participation:
\begin{definition}\label{def:bparticipation}
The \textbf{\bparticipation} schema is given by
\[
  \Pi_b = \left\{\pi \subseteq [n]  \mid \{i, j\} \subseteq \pi, i \ne j \Rightarrow |i - j| \ge \minsep\right\}.
\]
\end{definition}


Note that for \multiepoch-participation, $|\Pi_{(\maxpart,\minsep)}| = \minsep$, a fact \citet[Eq. 3]{choquette22multiepoch} critically exploited when computing sensitivity as brute force computation of $\max_{\pi \in \Pi}$.  With \bparticipation, we have $|\Pi_b| = \calO(\exp(n))$,  \btodo{Could compute something more exact here.} and hence such brute force approaches will be impractical even when we require $\bfC\tp \bfC \ge 0$.

Fortunately, this challenge can be surmounted for $\bands$-banded $\bfC$. We begin with the following Corollary.

\begin{cor}\label{cor:matrix_to_vector_sens}
For any dimensionality $\mdim \ge 1$, when per-step contributions are bounded by $\clipnorm=1$, consider a participation schema $\Pi \subseteq \Pi_b$ and matrix $\bfC$ such that for all $\pi \in \Pi$, for any $i, j \in \pi$ with $i \ne j$, we have $\left<\bfC\idx{:}{i}, \bfC\idx{:}{j} \right> = 0$ (that is, any pair of columns of $\bfC$ corresponding to distinct iterations $i, j$ where the same example might participate must be orthogonal).  Then,
\[
\senss(\bfC) = \sensv(\bfC).
\]
\end{cor}
\begin{proof}
The orthogonality condition implies $\bfC\idx{:}{\pi}\tp \bfC\idx{:}{\pi}$ is a diagonal matrix, with the (non-negative) norms of the selected columns forming the diagonal, and hence Condition 2 of \citet[Thm. G.1]{choquette22multiepoch} applies to $\bfC\idx{:}{\pi}$, and the result follows.
\end{proof}

Observe that a sufficient condition for \cref{cor:matrix_to_vector_sens} to hold is that $\bfC$ is $\minsep$-banded.\footnote{A stronger version of \cref{cor:matrix_to_vector_sens} can of course be proved immediately, requiring only  $\left<\bfC\idx{:}{i}, \bfC\idx{:}{j} \right> \ge 0$. However, this extra generality is not necessary for our purposes.} In the following theorem, we present two practical approaches to bounding sensitivity:


\newcommand{\vals}{c}  % Values being optimized
% TODO - consider switching i and t in the code
% and here ... to be t is better as a "budget" and
% i is better as an index.
\newcommand{\bi}{i}  % Total number of elements to choose, including x(t)
\newcommand{\ridx}{t} % Rightmost index
\algnewcommand{\LineComment}[1]{\State \(\triangleright\) #1}


\begin{algorithm}[t]
\caption{Efficient sensitivity calculation for \bparticipation} \label{alg:minsepsens}
\begin{algorithmic}%[1]
\State \textbf{Inputs:} min-separation $\minsep$, Gram matrix $\bfX$, max participations $\maxpart$
\State Initialize $F \in \R^{n \times k}$
\For{$m = 1, \dots, k$}
\For{$i = n, \dots, 1$}
 \Comment{We use the convention that $F[s,t] = 0$ if $s,t$ are out-of-bounds.}
\State $F[i, m] = \max{\Big(\bfX_{ii} + F[i+b, m-1], F[i+1, m]\Big)}$
\EndFor
\EndFor
\Return $\sqrt{F[1, k]}$
\end{algorithmic}
\end{algorithm}

\newcommand{\indpi}{u(\pi)}

\begin{theorem}\label{thm:minsepsens}
Let $\bfC \in \R^\dimdim$ be a lower-triangular $\bands$-banded matrix, and and the \bparticipation schema $\Pi_b$.  Further, suppose $k'$ upper-bounds the actual maximum number of participations that occurred in a data stream $\obs$ (at worst, we may take $k' \le \ceil{\frac{n}{b}}$):
Then:
\begin{enumerate}
\item \cref{alg:minsepsens} computes $\senss(\bfC)$ under schema $\Pi_b$ in polynomial time, and if $k' < \ceil{\frac{n}{b}}$, may produce a tighter bound on sensitivity for streams matching the actual participations in $\obs$.\footnote{Recall that for our notion of adjacency, all adjacent streams $\obsp$ must have this property.}
\item If $\kappa$ is an upper-bound on the column norms of $\bfC$, that is $\ \forall j \in [n], \ \norm{\bfC\idx{:}{j}} \le \kappa$, then the sensitivity is bounded by $\kappa \sqrt{k'}$.
\end{enumerate}
\end{theorem}
%
\begin{proof}
For \cref{alg:minsepsens}, let $\vals \in \R^\dimension$ with entries $\vals_i = \norm{\bfC\idx{:}{i}}^2$ for $i \in \{0, \dots, \dimension-1\}$. We have
\begin{equation}\label{eq:orthcolnorm}
\senss(\bfC) 
  = \max_{\pi \in \Pi_b} \norm{\bfC \indpi} 
  = \max_{\pi \in \Pi_b} \norm{\sum_{i \in \pi} \bfC\idx{:}{i}}  
  = \max_{\pi \in \Pi_b} \sqrt{\sum_{i \in \pi} c_i}
\end{equation}
where $\indpi \in \{0, 1\}^\dimension$ is given by $\indpi_i = 1$ if $i \in \pi$ and $0$ otherwise. The last equality follows from the orthogonality condition on sufficiently separated columns of $\bfC$. It is straightforward to verify the dynamic program of \cref{alg:minsepsens} constructs a feasible $\pi$ which attains the maximum.
\end{proof}

We find that $\bfC$ matrices normalized to have equal column norms perform well, and so hence often the $\kappa \sqrt{\minsep}$ bound of \cref{thm:minsepsens} is tight, and can easily be incorporated into optimization (in fact, the constant scaling can be handled post-optimization, so it is sufficient to simply enforce a fixed maximum column norm).